\cgeentry{cge_nash}{Nash equilibria of players who choose instruments, and leaders}

\cgesec{Description}

\cmd{cge_nash} (from version 0.20.0) finds the Nash equilibrium of players who
each choose instruments to maximise or minimise an outcome of the model,
given the others' choices, as an equilibrium problem in GAMS's EMP framework
does, and the choice of a leader who anticipates the others' equilibrium, as
in a Stackelberg game. A player is defined with \cmd{cge_nash player}, with
its instruments, objective and bounds written as in \manref{cge_optimize};
\cmd{cge_nash solve} finds the equilibrium and stores it as a simulation.

\cgesec{Quick start}

Two firms that each choose an output to maximise profit, in a standalone
model where the price depends on both outputs
\begin{stcode}
cge_nash player firm1, instruments(Q1) maximize((P - c) * Q1) lower(Q1:0)
cge_nash player firm2, instruments(Q2) maximize((P - c) * Q2) lower(Q2:0)
cge_nash solve, name(cournot)
\end{stcode}
Firm 1 as the leader
\begin{stcode}
cge_nash solve, name(stackelberg) leader(firm1)
\end{stcode}

\cgesec{Syntax}

\begin{stsyntax}
\cmdb{cge_nash player} \argu{name} \texttt{,} \optarg{inst}{ruments}{list} \{\optarg{max}{imize}{objective} $|$ \optarg{min}{imize}{objective}\} [\optarg{low}{er}{spec} \optarg{up}{per}{spec}]\\
\cmdb{cge_nash solve} \texttt{,} \optarg{}{name}{name} [\optarg{lead}{er}{player} \optarg{tol}{erance}{\#} \optarg{iter}{ate}{\#} \opt{}{trace}]\\
\cmdb{cge_nash list}\\
\cmdb{cge_nash clear} [\argu{player}]
\end{stsyntax}

\begin{opttable}
\optgroup{Players}
\optarg{inst}{ruments}{list} & the player's instruments, as in \manref{cge_optimize}\\
\optarg{max}{imize}{objective} & a report measure or expression to maximise\\
\optarg{min}{imize}{objective} & a report measure or expression to minimise\\
\optarg{low}{er}{spec}, \optarg{up}{per}{spec} & bounds, as \argu{instrument}\cmd{:}\argu{value}\\
\optgroup{Solve}
\optarg{}{name}{name} & the simulation of the equilibrium\\
\optarg{lead}{er}{player} & the leader, who anticipates the others' equilibrium\\
\optarg{tol}{erance}{\#} & largest change of an instrument in a round; default \texttt{1e-8}\\
\optarg{iter}{ate}{\#} & rounds of best responses; default 50\\
\opt{}{trace} & display each round\\
\end{opttable}

\cgesec{Options}

\begin{optdesc}
\item[\optarg{inst}{ruments}{list}] gives the player's instruments, as in
\manref{cge_optimize}: elements of parameters that \manref{cge_shock} changes,
of user parameters or of variables the closure fixes. An instrument belongs to
one player.

\item[\optarg{max}{imize}{objective}] and \optarg{min}{imize}{objective} give the
player's objective, a report measure or an expression of the user language.

\item[\optarg{low}{er}{spec}] and \optarg{up}{per}{spec} bound the
instruments. An instrument at a bound, where a best response may leave it,
starts the next one just inside it.

\item[\optarg{lead}{er}{player}] makes \argu{player} a leader: its instruments
are chosen to maximise its objective at the equilibrium of the other players
that each choice gives.

\item[\optarg{tol}{erance}{\#}] ends the rounds of best responses when none
changes an instrument by more than \# relative to its size.
\end{optdesc}

\cgesec{Remarks and examples}

In each round, each player in turn chooses its best response to the others'
instruments, by \manref{cge_optimize} with exact derivatives, and leaves it in
place for the next (a Gauss--Seidel iteration). At a Nash equilibrium no
player can do better given the others' choices, and a round changes nothing.
The iteration converges when the players' responses are not too strongly
interdependent, as in the examples; it may cycle otherwise.

\cgesubsec{Example 1: Cournot and Stackelberg}

Inverse demand $P = a - b(q_1 + q_2)$ with $a = 10$, $b = 1$ and a unit cost
$c = 1$; the Cournot outputs are $(a - c)/3b = 3$, and a leader produces $(a -
c)/2b = 4.5$, against the follower's $(a - c)/4b = 2.25$:
\stlog{nash_cournot}
The leader's profit, $(a - c)^2/8b = 10.125$, exceeds its Cournot profit of 9.

\cgesec{Stored results}

\cmd{cge_nash solve} stores the following in \cmd{r()}:
\begin{storedres}
\resgroup{Scalars}
\res{r(converged)} & 1 if the equilibrium was found\\
\res{r(rounds)} & rounds of best responses (without a leader)\\
\res{r(N_players)} & players\\
\res{r(objective)} & the leader's objective (with a leader)\\
\res{r(N_equilibria)} & the followers' equilibria solved (with a leader)\\
\resgroup{Macros}
\res{r(name)} & the simulation of the equilibrium\\
\res{r(leader)} & the leader, if any\\
\resgroup{Matrices}
\res{r(instruments)} & the players' instruments at the equilibrium\\
\res{r(objectives)} & their objectives\\
\res{r(leader_instruments)} & the leader's instruments (with a leader)\\
\end{storedres}

\cgesec{Methods and formulas}

With players $p = 1, \dots, P$, instruments $u_p$ and objectives $f_p(x(u),
u)$, a Nash equilibrium $u^*$ satisfies $u_p^* \in \arg\max_{u_p} f_p(x(u_p,
u_{-p}^*), u_p, u_{-p}^*)$ for every $p$. Each round sets $u_p \leftarrow
\arg\max f_p(\cdot, u_{-p})$ for $p = 1, \dots, P$ in turn, until the largest
change $|\Delta u| / \max(1, |u|)$ of a round is below the tolerance. A leader
$\ell$ chooses $u_\ell$ to maximise $F(u_\ell) = f_\ell(x, u_\ell, u_{-\ell}^*(u_\ell))$,
with $u_{-\ell}^*(u_\ell)$ the followers' equilibrium, by a projected Newton
method with central differences of $F$ (steps $10^{-3}$ times the scale of
each instrument), each evaluation solving the followers' equilibrium.

\cgesec{Also see}

\begin{alsosee}
\seeentry{cge_optimize}{Choose policy instruments to maximise an outcome or reach targets}
\seeentry{cge_model}{A standalone model, written whole in the user language}
\end{alsosee}
